% ---------------------------------------------------------------
\section{Introduction}\label{sec:intro}
% ---------------------------------------------------------------

Every large language model is the output of a data pipeline. Web crawls, licensed corpora,
human annotations and, more and more, synthetic data from other models are filtered, mixed,
recorded and fed to a training job. Each of these stages is a point where an adversary can
act. Poisoning web-scale training sets is practical~\cite{carlini2024poisoning}. Small
poisoned fractions persist through pre-training into aligned models~\cite{zhang2025persistent}.
Outsourced training can quietly skip or alter steps~\cite{jia2021pol,spoczynski2025atlas}.
In response, a large body of work proposes \emph{provenance} for machine learning: signed
data manifests~\cite{stokes2021prevent}, attested pipelines and transparency
logs~\cite{spoczynski2025atlas,duddu2025laminator}, hash-chained audit
trails~\cite{li2025auditablellm}, blockchain-anchored records~\cite{geren2025blockchain,mehammed2025fl},
and proofs of training~\cite{jia2021pol,abbaszadeh2024zkpot,liao2026verilora}.

These proposals rarely state what an auditor may conclude when the checks pass. Some
claims are too strong. Provenance is said to \emph{prevent} poisoning~\cite{stokes2021prevent},
but an authorised source that supplies poison passes every provenance check. Other claims
are left implicit. A tamper-evident log of a training run says nothing about whether the
released model came from that run unless the computation is bound to the record. Surveys
catalogue the techniques~\cite{geren2025blockchain,lyu2026survey,peng2026zkp}, but none
answers the auditor's question: \emph{given which mechanisms, against which adversary, does
an accepted record prove what?}

Three kinds of assurance are often conflated in this literature, and we keep them apart:
\begin{itemize}
\item \emph{Provenance}: who supplied an item, and what happened to it before training.
\item \emph{Computation verification}: whether the released model was produced from a stated
  record by a stated procedure.
\item \emph{Content assurance}: whether the data is free of poison.
\end{itemize}
The first two are properties of records and computations. The third is a property of data.

\paragraph{Research questions}
\begin{itemize}
\item[\textbf{RQ1}] Which combinations of provenance mechanisms are sufficient, and which are
  minimal, for each assurance goal against each adversary?
\item[\textbf{RQ2}] Which goals can no combination of mechanisms achieve, and what is the
  strongest claim that remains?
\item[\textbf{RQ3}] Do implementations of the mechanisms behave as the model predicts? How
  large are the quantities the model leaves free, such as grinding influence and detector
  recall, in practice?
\end{itemize}

\paragraph{This paper}
We answer these questions for training-data pipelines. Our approach has three parts.

\emph{(1) An assurance model.} We abstract the deployed techniques into seven mechanism
primitives: two kinds of computation binding, source witnessing, input witnessing,
attested capture, attested transform, and an attested content detector. We define seven
adversary capabilities, one per point of trust, including corrupt witnesses and broken
trusted hardware. We define six assurance goals, from ``the record explains the model's
inputs'' to ``no poison was used''.

\emph{(2) An exact characterisation.} A separability lemma shows that every instance of the
model can be decided by reasoning about a single item. We therefore decide all 9,216
(adversary, mechanism set, pipeline type, goal) instances for six named adversary profiles
exactly, with two independent procedures that agree on every one. We then extend the
analysis to all $2^7=128$ subsets of adversary capabilities (196,608 decisions); every
theorem holds over this full lattice. The minimal sufficient mechanism sets
(Table~\ref{tab:minimal}) yield six theorems:
\begin{itemize}
\item Computation binding is necessary exactly when the training executor is untrusted.
\item Lineage has two routes, source witnessing or attested capture-plus-transform, which
  rest on disjoint trust anchors.
\item Generative pipelines need declared and authenticated stage types, or injected data can
  be \emph{laundered} as synthetic.
\item A binding check that accepts a prover-chosen seed leaks influence. We bound it at
  $\log_2T$ bits for $T$ training runs.
\item No mechanism set certifies the absence of poison. The strongest available claim is
  conditional on a detector's recall.
\item Record completeness can be achieved by attestation or by witnessing, which fail under
  disjoint assumptions. This is a formal argument for trust diversity.
\end{itemize}
A ledger adds censorship-resistant inclusion to a witnessed transparency log, and nothing
else that the goals depend on.

\emph{(3) Empirical validation.} The theorems concern a model, so we test whether the model
is faithful. We built a prototype pipeline with real cryptography (Ed25519 receipts, Merkle
commitments, a witnessed transparency log), trusted execution \emph{simulated} by key
isolation, real training on a public sentiment corpus with
exact re-execution, a published backdoor detector~\cite{tran2018spectral}, and an adaptive
adversary that tries every combination of its available attacks. Across all six adversary
profiles and 128 mechanism sets, no execution that the model says is safe violated a goal,
and the attack library realised the model's predicted failures (Section~\ref{sec:eval}). A
grinding experiment with 1,024 training runs shows that a prover can set
$\lfloor\log_2T\rfloor$ chosen bits of model behaviour while passing a re-execution check,
which matches Theorem~\ref{thm:grind} to within one bit. On a single consumer GPU, the same experiments with DistilBERT and Qwen2.5-0.5B~+~LoRA give
three results. Fine-tuning replays bit for bit on fixed hardware. The grinding curve is
identical: 7 chosen bits for 128 DistilBERT runs. A backdoor delivered at 2\% by an authorised source passes every provenance check.
Across three seeds, it raised attack success by 20 to 75 points over the clean model, while
spectral-signature recall averaged 25\%. These GPU results cover fine-tuning with one model
family and one detector; they do not cover pre-training.

\paragraph{Contributions}
\begin{itemize}
\item A systematisation of training-data provenance mechanisms as seven primitives, with a
  mapping of published systems to them (Section~\ref{sec:sok}).
\item A formal assurance model with six goals and seven adversary capabilities, together with
  a separability lemma that makes it exactly decidable (Sections~\ref{sec:model}
  and~\ref{sec:results}).
\item A complete characterisation of minimal sufficient mechanism sets for six named
  adversary profiles; six theorems with proofs; and a check of every theorem over all 128
  adversary capability subsets (Section~\ref{sec:results}, Appendix~\ref{app:proofs}).
\item A prototype and evaluation that test the model's soundness and tightness on real
  executions, a measurement of seed grinding against the bound, and overhead figures
  (Section~\ref{sec:eval}).
\item Design guidance and open problems for auditors, system builders and regulators
  (Section~\ref{sec:discussion}).
\end{itemize}
All models, scripts, raw results and screening logs form an artefact that will be deposited
for review (Section~\ref{sec:artefact}).
